\section{Schuon vs Guenon}

\begin{quotex}
The true essence of everything always abides, though unmanifest in the inner depths of very Being, while its sensory qualities appear outwardly. For it is impossible that the Divine Ideas included in the intelligible world should be evanescent: to pretend that the content of Divine Science is evanescent would imply atheism.
\flright{\textsc{Nur ad-Din Abd ar-Rahman Jami}, \emph{Lawa'ih}}
\end{quotex}


The descent of the Being, or the Self, to the human being, ends with the subtle and gross states. The subtle state has three sheaths, as shown in the following chart.

\begin{itemize}


\item Corporeal things

  \begin{itemize}

  
\item Living beings, plants, vegetative soul

    \begin{itemize}

    
\item Sentient beings, animals, animal soul

      \begin{itemize}

      
\item Rational beings, humans, rational soul
      \end{itemize}

    \end{itemize}

  \end{itemize}

\end{itemize}


Of course, the descent happens in reverse order, from the rational to the gross; the chart illustrates how it unfolds in time, as material objects seem to evolve from life to sentience to rationality. The subtle state may also be called the psychical. Note that the gross state and the subtle state together make up what may be called the corporeal world. That is why the following quote from \textbf{Frithjof Schuon} seems to be way off the mark:

\begin{quotex}
I do not know from where Guenon gets this enumeration of the five conditions of physical existence which he calls `corporeal': space, time. form, number, and life. l am in agreement as regards the first four, but not as regards life, because what we are concerned with here is matter or, if one prefers, physical substance. If one adds life, which is not at all a general condition, it is likewise necessary to add other secondary categories such as color and so on.
\flright{\textsc{Frithjof Schuon}, \emph{Rene Guenon: Some Observations}}
\end{quotex}


\paragraph{Five Conditions}


\emph{Pace} Schuon, \textbf{Rene Guenon} considers this to be a very important question, so it is important to ``get it right''. These are Guenon's actual words, which Schuon misquoted or mistranslated:

\begin{quotex}
The existence of individual beings in the physical world is indeed subject to a set of five conditions: space, time, matter, form and life, which can be matched with the five bodily senses, as well as to the five elements.
\flright{\textsc{Rene Guenon}, \emph{Gnosis and Spiritualist Schools}}
\end{quotex}


Guenon identifies the source as the \emph{Sankhya} of Kapila. The \emph{tanmatras} are the five elements, which are related to the five senses, as shown in \emph{Man and his Becoming}. Likewise the five elements are related to those elements. Therefore, the chief characteristic of the five conditions is their transcendence in the sense that they are never themselves manifested, yet are the principles of manifestation.

\textbf{Space} and \textbf{Time} are obviously like that; they are \emph{a priori} to experience. That is, there cannot be any phenomenal experience without space and time. Space is not discovered by investigating the world, as if there could be someplace that is non-space. The same applies to time.

There is a further distinction. Time arises in the subtle state, but space arises in the gross state. Ether is the most subtle element, at the border between the gross state and the etheric body of the subtle state.

By ``\textbf{form}'', Guenon means the ideas or essences in the Divine Intellect. Once again, they are not discovered through investigation of the world. Forms are not grasped by the senses, but rather by the intellect intuitively.

There is another possible misunderstanding: Guenon uses the terms ``physical world'' and ``corporeal world'' interchangeably. Hence, the corporeal world includes the subtle world. Perhaps the Greek word \textbf{psyche} could be used instead of ``\textbf{life}'' to make the point more clearly. Psyche is ``life'', but it has connotations like soul or consciousness. Now consciousness is indeed a ``general condition'', irreducible to the other conditions; consciousness does not arise out of matter. Nor is consciousness discovered in the phenomenal world; on the contrary there is no phenomenal world without consciousness.

Colours can only arise in consciousness so they are indeed secondary. This dissolves Schuon's objections. \textbf{Matter} will be taken up in the following section.

\paragraph{The Notion of Matter}


After having addressed the issue of ``life'' as general principle, Schuon turns to \textbf{matter}:

\begin{quotex}
For Guenon the notion of `matter', is factitious, confused. problematic; it has nothing fundamental about it and is to be found nowhere except in the modern West. This is incredible. And what, in an altogether general way, is the sensible substance that one can touch, measure, weigh, analyze, and possibly work or shape? And why, for goodness' sake, would this not be matter?
\end{quotex}


Now even a great intellect can be wrong, but his errors are usually deep and subtle. If the mistake seems so obvious, then it is likely that something essential has been lost. What follows are a few points to take into consideration.

\textbf{Point 1: Matter is not manifested}


This means that you cannot show me a piece of matter. Like space and time, matter is part of each experience, yet it cannot be experienced directly. Hence, one cannot touch, measure, weigh, analyze, and possibly work or shape it directly.

\textbf{Point 2: Experience is of things}


Following up on point 1, you can show me a banana, an automobile, a squirrel, and so on, but you cannot show me matter. Hence matter is not what can touch, measure, weigh, analyze, and possibly work or shape. Rather, those experiences are of specific things. You can touch a banana, weigh it, analyze it, and slice it up over your ice cream.

\textbf{Point 3: Matter has no meaning in physics}


This may surprise some people, but the notion of matter is not fundamental to physics. You could define matter as any elementary particle with mass. Yet quantum entities do not have size or volume as conventionally understood. That means that the classical definition of matter as ``extension'' is not valid. However, corporeal things do have extension.

\textbf{Point 4: Matter is not necessary for sense experience}


You can experience sights, colours, sounds, the passage of time, travel, and so on in your dreams or imagination. Hence, sensory experience does not prove the existence of ``matter''.

\textbf{Point 5: Matter is indeterminate}


The ``matter'' of physics is more like the prime matter of ancient philosophy. That is, it is in an indeterminate state, just like Schrodinger's famous cat. It only becomes a specific thing when it is informed. Wolfgang Smith has addressed this question adequately, so there is no need to repeat it here.

\paragraph{Epilogue}


True understanding requires a shift in consciousness. The naïve view is that matter is reality and there are things ``out there, right now'' that run on their own power.

The shift is to understand matter as the canvas on which the Divine Archetypes are reflected. They manifest in space and time, so that they can be experienced in consciousness. It is a gift. You can grasp the Attributes in the intellect, but you can also experience them through the senses.


\flrightit{Posted on 2021-07-28 by Cologero}

\begin{center}* * *\end{center}

\begin{footnotesize}
\begin{sffamily}
\texttt{Presbyter Iohannes on 2021-07-29 at 13:37 said:}

What do you have to say about ``Schuon vs Guénon'' in more general terms?

\hfill

\texttt{Cologero on 2021-08-01 at 21:08 said:}

The most striking difference, in my opinion, is in their style. Guenon is more adamant, opposing one view to its contraries. Schuon strives to seek out the middle ground between two seemingly opposing views, in which each view has plausibility in its perspective. So Schuon addresses exoteric issues, which Guenon typically avoids.

\hfill

\texttt{Presbyter Iohannes on 2021-08-01 at 21:33 said:}

\url{https://pt.scribd.com/document/348127488/Rfutation-Des-Quelques-Critiques-de-Schuon}


This anonymous refutation covers all of the Schuon fiasco against Guénon. Unfortunately it is only available in French. Perhaps it would be useful if someone translated it (someone more well versed in french than I am).

\hfill

\texttt{Maude Murray on 2022-07-28 at 23:47 said:}

For an easy to understand comparison, see Frithjof Schuon versus René Guénon: their roles near the end of the world, by Maude Murray, the third of Schuon's four ``wives!'' That's on Academia.edu. Or see Twitter, Facebook or A last minute lesson in discernment, under her name.

\hfill

\texttt{Maude Murray on 2022-11-20 at 22:39 said:}

Just go to \url{https://frithjofschuon.wordpress.com} where the two are decisively compared in their own words. They are also compared to all Gnostics in all revelations from time immemorial. The difference is astounding!

I wish I could delete my past comments here as some sites have been deleted or changed (hint).

\texttt{Kaukomieli on 2022-12-18 at 08:49 said:}

Here is a short video on YouTube where Schuon discusses his views on metaphysics, esotericism, syncretism and religion.

\url{https://youtu.be/OKgjGiTVbPo}


Two striking differences between them seems to be in their relationship to syncretism and exoteric religious practice. Where as Guenon rejects syncretism and speaks of synthesis, Schuon seems to accept syncretism in esoteric practice, and whereas Guenon stresses the need for religious exotericism, Schuon seems to considered exoteric religious practice unnecessary for an esotericist or to a metaphysician.

I've also heard Schuon was also quite extravagant which derived from his experience of Mary from which he made his practice of ritual nudity going all the way to the practice of sexual magic. Some have said it was a sign of a spiritual illness rather than a mystical revelation.

\hfill

\texttt{Kaukomieli on 2022-12-18 at 13:50 said:}

Here is also a lenghty document of the two gentlemen in question that might be of interest:

\url{https://youtu.be/7Vlf8GQJFeY}


Today I read some texts that describes Schuon as defiling minors with his cultic ritual, but what I gathered was that the profane world had misunderstood his apparently illegal activities with the Girls in their teens. For the ritual setting of the Khrishnaic dance with its touching of nude human bodies is completely out of the moral framework of sexual license.

\hfill

\end{sffamily}
\end{footnotesize}

